\documentclass[11pt]{article}
\usepackage[margin=1in]{geometry}
\usepackage{amsmath,amssymb,amsthm}
\usepackage{booktabs}
\usepackage{microtype}
\usepackage[svgnames]{xcolor}
\usepackage[colorlinks=true,linkcolor=Teal,citecolor=Teal,urlcolor=Teal]{hyperref}

\newcommand{\R}{\mathbb{R}}
\newcommand{\C}{\mathbb{C}}
\newcommand{\hd}{d_{\mathrm{head}}}
\newcommand{\nb}{n_b}
\DeclareMathOperator{\softplus}{softplus}
\DeclareMathOperator{\atantwo}{atan2}
\DeclareMathOperator*{\cumsum}{cumsum}

\title{Path Integration and PoPE:\\ how the two mechanisms compose}
\author{Notes derived from the implementation}
\date{}

\begin{document}
\maketitle

\begin{abstract}
\noindent
Every positional mechanism here is a choice about one of two factors in the same
product. Writing an attention score in polar form as $\mu\,e^{i\varphi}$,
\textbf{MapFormer's path integration determines the positional part of the phase
$\varphi$} and \textbf{PoPE determines the magnitude $\mu$}, which is why the
combination (\textsc{MapPoPE}) is not a new mechanism but the conjunction of two
independent choices.

\medskip
\noindent\textbf{The symmetry is not exact, and the asymmetry matters.} MapFormer
\emph{adds} a term to the phase and leaves the rest of the attention untouched.
PoPE does two things: it replaces the magnitude \emph{and} it eliminates the
content-dependent intrinsic phase $\psi$ that ordinary rotary attention carries.
Under PoPE each element is a non-negative real times a purely positional phasor, so
$\tilde q$ can no longer lie anywhere in $\C$. PoPE is therefore a
\emph{restriction} on the object MapFormer rotates, not merely a different
parameterisation of one factor --- and that is a plausible source of the
sub-additivity observed on text ($\S$6).

\medskip
\noindent Formulas below are transcribed from the implementation, and where the
implementation departs from the published equation or from canonical RoPE that is
stated at the point of departure.
\end{abstract}

\section{The shared frame}

Let $x_1,\dots,x_T \in \R^{d}$ be token embeddings, and for a single head write
queries and keys $q_t = W_Q x_t$, $k_s = W_K x_s \in \R^{\hd}$.
Group the $\hd$ coordinates into $\nb$ two-dimensional blocks and identify each
block with a complex number,
\begin{equation}
  \tilde q_t^{(c)} \;=\; q_{t,2c} + i\,q_{t,2c+1},
  \qquad
  \tilde k_s^{(c)} \;=\; k_{s,2c} + i\,k_{s,2c+1},
  \qquad c = 0,\dots,\nb-1 .
\end{equation}

A rotary scheme attaches a per-position angle $\theta_t^{(c)}$ and rotates:
\begin{equation}
  \label{eq:rot}
  \begin{pmatrix} q'_{t,2c} \\ q'_{t,2c+1}\end{pmatrix}
  =
  \underbrace{\begin{pmatrix}
    \cos\theta_t^{(c)} & -\sin\theta_t^{(c)}\\
    \sin\theta_t^{(c)} & \phantom{-}\cos\theta_t^{(c)}
  \end{pmatrix}}_{R(\theta_t^{(c)})\,\in\,SO(2)}
  \begin{pmatrix} q_{t,2c}\\ q_{t,2c+1}\end{pmatrix},
\end{equation}
and likewise for $k_s$. The unnormalised attention logit is then
\begin{equation}
  \label{eq:score}
  a_{ts}
  \;=\;\frac{1}{\sqrt{\hd}}\,\langle q'_t,\, k'_s\rangle
  \;=\;\frac{1}{\sqrt{\hd}}\sum_{c=0}^{\nb-1}
        \big\langle R(\theta_t^{(c)})\,q_t^{(c)},\; R(\theta_s^{(c)})\,k_s^{(c)}\big\rangle .
\end{equation}

Because $R$ is orthogonal, $R(\alpha)^{\!\top} R(\beta) = R(\beta-\alpha)$, so
\begin{equation}
  \label{eq:relative}
  a_{ts}
  \;=\;\frac{1}{\sqrt{\hd}}\sum_c
       \big\langle q_t^{(c)},\, R\!\left(\theta_s^{(c)}-\theta_t^{(c)}\right) k_s^{(c)}\big\rangle
  \;=\;\frac{1}{\sqrt{\hd}}\sum_c
       \operatorname{Re}\!\Big[\overline{\tilde q_t^{(c)}}\;\tilde k_s^{(c)}\;
         e^{\,i(\theta_s^{(c)}-\theta_t^{(c)})}\Big].
\end{equation}
Only the \emph{difference} of angles survives: this is the relative-position
property, and it holds for \emph{any} choice of $\theta$.

Finally, write each complex factor in polar form,
$\tilde q_t^{(c)} = \mu^q_{t,c}\,e^{i\psi^q_{t,c}}$ and similarly for $k$.
Then \eqref{eq:relative} becomes
\begin{equation}
  \label{eq:polar}
  \boxed{\;
  a_{ts} \;=\; \frac{1}{\sqrt{\hd}}\sum_c
     \underbrace{\mu^q_{t,c}\,\mu^k_{s,c}}_{\text{magnitude}}\;
     \cos\!\big(\underbrace{\theta_s^{(c)}-\theta_t^{(c)} + \psi^k_{s,c}-\psi^q_{t,c}}_{\text{phase}}\big).}
\end{equation}
Equation~\eqref{eq:polar} is the frame for everything that follows. \textbf{Each
mechanism is a decision about where the magnitude comes from and where the phase
comes from.}

\section{RoPE: the fixed clock}

Standard RoPE fixes the angle to be linear in the token index,
\begin{equation}
  \theta_t^{(c)} \;=\; t\cdot\omega_c,
  \qquad
  \omega_c \;=\; \text{base}^{-2c/\hd} \;=\; \text{base}^{-c/\nb},
\end{equation}
with $\omega$ a non-learnable buffer. The magnitude and the intrinsic phase
$\psi$ are whatever $W_Q, W_K$ happen to produce; they are never constrained, and
in particular $\tilde q$ may lie anywhere in $\C$. Substituting into
\eqref{eq:relative} gives the familiar dependence on $s-t$ alone.

\section{MapFormer: path integration replaces the clock}

MapFormer keeps \eqref{eq:rot} and changes only what generates $\theta$.

\paragraph{Step 1 --- token to Lie-algebra increment.}
A low-rank map (rank $r=2$; paper App.~A.7, $W_\Delta = W_\Delta^{\mathrm{out}}
W_\Delta^{\mathrm{in}}$) reads \emph{every} token, action or observation:
\begin{equation}
  \Delta_t \;=\; W_{\mathrm{out}}\,W_{\mathrm{in}}\,x_t \;\in\;\R^{H\times \nb},
  \qquad
  W_{\mathrm{in}}\in\R^{r\times d},\quad W_{\mathrm{out}}\in\R^{(H\nb)\times r}.
\end{equation}
Nothing tells the model which tokens are actions; it must \emph{learn}
$\Delta \approx 0$ on observations. The paper states only $r \ll d$, ``for instance
$r=2$'' (App.~A.7).

\begin{quote}\small
\textbf{Unverified.} A note in this project's log records the bottleneck as
load-bearing, with $r=1$ dropping the 2D task to $0.66$. \emph{No experiment in
this repository varies $r$ on the 2D task}, and the locally held paper version
(v3) contains no rank ablation. The claim is repeated here only to be marked
unsupported; do not cite it.
\end{quote}

\paragraph{Step 2 --- integrate.}
\begin{equation}
  \label{eq:pi}
  \boxed{\;\theta_t^{(h,c)} \;=\; \omega^{(h,c)}\sum_{u\le t}\Delta_u^{(h,c)}
  \;=\;\omega^{(h,c)}\,\cumsum(\Delta)^{(h,c)}_t . }
\end{equation}
Because $SO(2)$ is abelian, the prefix product of rotations collapses to a
\emph{cumulative sum of angles}, which is what makes this an $O(\log T)$ parallel
scan rather than a sequential recurrence:
\begin{equation}
  \prod_{u\le t} R\!\left(\omega\Delta_u\right) \;=\; R\!\Big(\omega\textstyle\sum_{u\le t}\Delta_u\Big).
\end{equation}
Commutativity is the entire computational argument, and it is also the limitation:
a turn/forward action space, where displacement depends on accumulated heading, is
\emph{not} representable by \eqref{eq:pi}.

\paragraph{Step 3 --- frequency schedule.}
The $\omega$ are learnable, initialised geometrically over $\nb$ blocks:
\begin{equation}
  \label{eq:omega}
  \omega_i \;=\; \omega_{\max}\left(\frac{1}{\Delta_{\max}}\right)^{\frac{i}{\nb-1}},
  \qquad i = 0,\dots,\nb-1,
  \qquad \omega_{\max}=2\pi,\quad \Delta_{\max}=\text{grid size}.
\end{equation}

\begin{quote}\small
\textbf{Correction to the published equation.} Eq.~\eqref{eq:omega} differs from
the paper (v1/v3 eq.~17, v4 eq.~18) in \emph{two} places, both deliberate. First
the \textbf{sign}: the paper prints a negative exponent, giving an \emph{increasing}
schedule that contradicts its own stated boundary condition $\omega_{\min} =
\omega_{\max}/\Delta_{\max} \le \omega_i \le \omega_{\max}$. Second the
\textbf{denominator}: the paper writes $n_b$ with $1 \le i \le n_b$, so the endpoint
$\omega_{\min}$ is never attained; eq.~\eqref{eq:omega} uses $n_b-1$ with
$0 \le i \le n_b-1$, which hits both endpoints exactly. The corrected form is what
satisfies the paper's own inequality and what matches RoPE's
$\theta_i = \text{base}^{-2i/d}$. Both defects survive into v4.
\end{quote}

\paragraph{RoPE is a special case --- of the functional form, not the encoding.}
If the model learns $\Delta_t \equiv 1$ for all $t$, \eqref{eq:pi} becomes
$\theta_t = (t{+}1)\,\omega$ (the cumsum is inclusive; only differences enter
\eqref{eq:relative}, so the offset is immaterial). That is RoPE's \emph{form}, with
the token index as the clock. It is \emph{not} RoPE identically, for three reasons:
$\omega$ here is a trained parameter rather than a frozen buffer; its schedule
spans $[2\pi/64,\,2\pi] = [0.098,\,6.28]$ against RoPE's $[10^{-4},\,1]$, a ratio
running from $6.3\times$ at $c=0$ to $982\times$ at $c=\nb-1$; and $\Delta_t$ is a
rank-2 function of the embedding, so $\Delta\equiv 1$ is a hypothesis about what
was learned, not a setting. The accurate statement is that path integration is
\textbf{RoPE with a learned, content-dependent clock rate --- which may be zero or
negative --- over a learned frequency spectrum}.

\medskip
\noindent In terms of \eqref{eq:polar}: MapFormer sets the \textbf{phase} and
leaves the magnitude untouched.

\section{PoPE: an explicit polar parameterisation}

PoPE (Polar Coordinate Positional Embeddings) makes \eqref{eq:polar} explicit
rather than emergent.

\begin{quote}\small
\textbf{Notation change, from here on.} $\S$1 defined $\nb$ as the number of
\emph{2D blocks}, $\nb = \hd/2$. PoPE indexes one frequency per \emph{element}, so
throughout this section and the next $\nb = \hd$. Read $\vartheta_c =
\text{base}^{-c/\nb}$ below as $\text{base}^{-c/\hd}$; under $\S$1's definition it
would be wrong by a factor of two in the exponent.
\end{quote}

\noindent Per element $c$:
\begin{align}
  \mu^q_{t,c} &= \softplus\!\big(Q_{t,c}\big), &
  \mu^k_{s,c} &= \softplus\!\big(K_{s,c}\big), &&\text{(content, strictly positive)}\\
  \varphi^q_{t,c} &= t\,\vartheta_c, &
  \varphi^k_{s,c} &= s\,\vartheta_c + \delta_c, &&\text{(position, plus a bias)}
\end{align}
with $\vartheta_c = \text{base}^{-c/\nb}$ a \textbf{fixed buffer} and
$\delta_c$ a learnable per-(head, frequency) constant.

\begin{quote}\small
\textbf{$\delta_c$ is very nearly inert in practice.} It is clamped to $[-2\pi, 0]$
in the \emph{forward pass}, not projected onto that range, so its gradient is
identically zero outside the interval --- and it is initialised at exactly $0$,
i.e.\ \emph{on} the upper boundary. Any step that would push it positive kills its
gradient permanently. Measured on trained checkpoints (4 heads $\times$ 64
frequencies $\times$ 4 layers), \textbf{91--98\% of $\delta_c$ entries end at
$\ge 0$}, clamped to zero with dead gradient; per-layer means run $-0.0003$ to
$-0.018$ rad. For roughly 95\% of frequencies, \eqref{eq:pope} is effectively
$\cos((s-t)\vartheta_c)$. Treat the bias as a nominal degree of freedom, not a
working one.
\end{quote}
The logit is
\begin{equation}
  \label{eq:pope}
  \boxed{\;a_{ts} \;=\; \frac{1}{\sqrt{\hd}}\sum_{c}
    \softplus(Q_{t,c})\;\softplus(K_{s,c})\;
    \cos\!\big((s-t)\,\vartheta_c + \delta_c\big).}
\end{equation}
Implemented via the angle-difference expansion, so no complex arithmetic is
needed:
\begin{align}
  \cos^K_{s,c} &= \cos\varphi_{s,c}\cos\delta_c - \sin\varphi_{s,c}\sin\delta_c, &
  \sin^K_{s,c} &= \sin\varphi_{s,c}\cos\delta_c + \cos\varphi_{s,c}\sin\delta_c,\\
  a_{ts} &= \tfrac{1}{\sqrt{\hd}}\Big[
      (\mu^q\!\cos\varphi)_t \cdot (\mu^k\!\cos^K)_s^{\top}
    + (\mu^q\!\sin\varphi)_t \cdot (\mu^k\!\sin^K)_s^{\top}\Big].
\end{align}

\paragraph{What PoPE does \emph{not} do.} Its phase is a function of position
only. The magnitude is content-dependent; the \emph{angle} is not, apart from the
constant bias $\delta_c$. This is worth stating because the name invites the
opposite reading: PoPE is \textbf{not} in the CoPE / Selective-RoPE family, which
learn content-dependent \emph{angles}. It modifies the other factor.

\medskip
\noindent In terms of \eqref{eq:polar}: PoPE sets the \textbf{magnitude} and
leaves the phase as a fixed clock.

\section{MapPoPE: the conjunction}

Since MapFormer determines $\varphi$ and PoPE determines $\mu$, combining them
requires no new mechanism --- only substituting one into the other. Replace
PoPE's positional phase $t\vartheta_c$ with the path-integrated angle
\eqref{eq:pi}:
\begin{equation}
  \label{eq:mappope}
  \boxed{\;
  a_{ts} \;=\; \frac{1}{\sqrt{\hd}}\sum_{c}
    \underbrace{\softplus(Q_{t,c})\,\softplus(K_{s,c})}_{\text{PoPE: magnitude from content}}
    \;\cos\Big(
      \underbrace{\theta_s^{(c)} - \theta_t^{(c)}}_{\text{MapFormer: phase from the path integral}}
      + \;\underbrace{\delta_c}_{\text{PoPE bias}}\Big),}
\end{equation}
with $\theta^{(c)}_t = \omega^{(c)}\sum_{u\le t}\Delta^{(c)}_u$ as in
\eqref{eq:pi}. In code the class is three lines: inherit MapFormer's forward pass
(so $\cos\theta,\sin\theta$ still come from the path integrator), widen the angle
machinery, and swap the attention module.

\paragraph{The one non-trivial detail: frequency count.}
RoPE rotates \emph{pairs} of coordinates, so it needs $\nb = \hd/2$ angles. PoPE
assigns one frequency per \emph{element}, so it needs $\nb = \hd$. The combination
therefore rebuilds both position modules at the wider count:
\begin{equation}
  \nb \;:\; \hd/2 \;\longmapsto\; \hd,
  \qquad
  W_{\mathrm{out}} \in \R^{(H\hd)\times r},
  \qquad
  \omega \in \R^{H\times \hd}.
\end{equation}
In parameters this is small: at $d_{\text{model}}=128$, $H=2$, one layer, the two
models are \textbf{204,373} and \textbf{204,693}, a difference of $+320$
($+0.16\%$) --- $+64$ on $\omega$, $+128$ on $W_{\mathrm{out}}$, $+128$ for
$\delta_c$. The consequential change is not the parameter count but the
\textbf{frequency count}: \textsc{MapPoPE} carries twice as many independent
position frequencies, indexed per element rather than per pair. So the two are
neither exactly parameter-matched nor frequency-matched, and the second is what a
comparison has to contend with.

\begin{quote}\small
\textbf{What this does and does not contaminate.} In the $2\times2$
$\{$RoPE, PoPE$\}\times\{$index, path-integrated$\}$, both PoPE cells are wide and
both RoPE cells narrow, so the width confound is \emph{constant along the encoding
factor}. The position main effect and the interaction are therefore clean; only the
encoding main effect carries it --- and that is the small one ($+0.011$ on the
torus against $+0.461$ for position). There is also a direct empirical check: if
the extra frequencies were doing the work, PoPE with an index phase should beat
RoPE with an index phase. It does not ($0.509$ vs $0.530$, both on the $0.506$
floor). The missing control is a \emph{narrow} PoPE at $\nb = \hd/2$ with each
angle repeated across its pair, which would isolate the parameterisation from the
width; it has not been run.
\end{quote}

\noindent The widening is a property of PoPE itself, not of the combination ---
the index-phase PoPE baseline also runs at $\nb = \hd$.

\section{Why the two compose, and when they do not}

\eqref{eq:mappope} predicts the observed pattern.

\begin{itemize}
\item \textbf{PoPE alone is inert where position carries the signal.} On the torus
  navigation task, PoPE with an index phase scores $0.509$ against a measured
  floor of $0.506$. A better magnitude parameterisation is worth nothing when the
  phase encodes no spatial information.
\item \textbf{PoPE plus path integration is the strongest configuration measured.}
  $0.994$ on the same task. Once the phase carries genuine position, sharpening the
  content channel pays.
\item \textbf{On text the picture does not survive its own metric.} Reading the
  \emph{final} checkpoint gives PoPE alone $-0.0053$ bpc, path integration alone
  $-0.0041$, sum $-0.0093$, both together $-0.0058$ --- apparently sub-additive.
  But the within-run checkpoint standard deviation is $0.003$--$0.007$, larger than
  every one of those quantities, which is why the source file specifies the
  \textbf{mean of the last five checkpoints} instead. On that metric the numbers
  are $-0.0041$, $-0.0006$, sum $-0.0047$, both $-0.0052$: path integration alone
  nearly vanishes and the combination is mildly \emph{super}-additive --- the
  opposite sign. Both single-axis arms are also $n=1$. \textbf{Nothing about
  additivity on text is established here in either direction}; the two readings are
  reported so the fragility is visible rather than resolved by choosing one.
\end{itemize}

\noindent The general statement: the two mechanisms are independent as
\emph{parameterisations} --- they touch disjoint factors of \eqref{eq:polar} ---
but that does not make their \emph{effects} independent. Whether they add depends
on whether the task has separate work for a phase and a magnitude to do.

\appendix
\section{The other mechanisms, in the same notation}

\paragraph{Weight-shared recursion (the ``loop'').}
With $B_\vartheta$ a single transformer block and $n$ the loop count,
\begin{equation}
  h^{(0)} = \text{Emb}(x), \qquad
  h^{(j+1)} = B_\vartheta\!\left(h^{(j)};\,\cos\theta,\sin\theta\right),\quad j=0,\dots,n-1 .
\end{equation}
$\theta$ is computed \emph{once} from the token embeddings and reused, matching
both parents. Parameter count equals one layer; compute equals $n$ layers.

Carrying a \emph{bounded, gated correction} to $\theta$ across iterations ---
$\theta \leftarrow \theta_0 + g\cdot\tanh(\text{refine}(h))$ with $g$ initialised
to zero --- is a different model, and it is null: the learned gate settles at mean
$|g| \approx 0.083$ with inconsistent sign across seeds, capping the correction at
$0.14$ rad against a $\theta$ spanning $\sim 2\pi T$. The gate was verified
escapable (gradient $1.9\times 10^{-3}$ at zero), so the optimiser declined rather
than failed to express it. A genuine per-iteration recomputation of
$\omega\cdot\mathrm{cumsum}(\Delta(h^{(j)}))$ has \emph{not} been run.

\paragraph{Level 1.5 invariant EKF (withdrawn; recorded for completeness).}
A correction applied to $\theta$ before the rotation:
\begin{align}
  R_t &= \exp\!\big(\mathrm{clamp}(f_R(x_t),-5,5)\big) & &\text{per-token measurement noise}\\
  z_t &= \pi\tanh\!\big(f_z(x_t)\big) & &\text{measurement in }[-\pi,\pi]\\
  \nu_t &= \atantwo\!\big(\sin(z_t-\theta_t),\,\cos(z_t-\theta_t)\big) & &\text{wrapped innovation}\\
  K_t &= \Pi/(\Pi + R_t),\qquad \Pi = e^{\log\Pi} & &\text{gain; }\Pi\in\R^{H\times \nb}\text{, constant in }t\\
  d_t &= (1-K_t)\,d_{t-1} + K_t\,\nu_t, \qquad d_{-1}=0 & &\text{affine scan}\\
  \hat\theta_t &= \theta_t + d_t . & &
\end{align}
$\Pi$ is a learned $(H\times\nb)$ constant --- 64 values at the default config, one
per head and frequency block --- not a scalar; it is constant across \emph{tokens},
which is what distinguishes Level 1.5 from the full heteroscedastic filter. The
recurrence is scalar in each channel and affine, so it is one more parallel scan.

\paragraph{Initialisation, which turned out to matter.} $\log R$'s head is
initialised with weights scaled by $0.01$ and bias $0$, giving $R\approx 1$ and
hence $K \approx 0.5$ \emph{at initialisation} --- corrections at half strength
before anything is learned. On the multiplicative-attention backbone this destroyed
the gradient path and diverged 3 of 9 seeds; the fix was a bias of $3.0$, i.e.\
$K\approx 0.05$, making the filter a near-no-op at init.

\paragraph{Status.} No individual component is load-bearing: the measurement head,
the per-token gate and the learned $\Pi$ can each be removed alone at no measurable
cost ($n=5$). The whole filter does have an effect, but only at out-of-distribution
length ($+0.062$ at $T{=}512$, $+0.124$ at $T{=}1024$, loss-matched), and on a task
built to supply the drift it is meant to correct the effect \emph{does not grow
with the drift} --- flat in one recipe, negative in another. The mechanism is
withdrawn as inference; what survives is stabilisation.

\paragraph{Hierarchy (hourglass).} With pooling factor $k$, a causal
\emph{shift-then-mean} pool $P_k$ and an upsample $U_k$,
\begin{equation}
  h \leftarrow B_{\mathrm{pre}}(h), \qquad
  h \leftarrow h + U_k\,B_{\mathrm{coarse}}\!\big(P_k h\big), \qquad
  h \leftarrow B_{\mathrm{post}}(h),
\end{equation}
with the coarse angle pooled alongside the content, and the sequence padded to a
multiple of $k$ then truncated after $B_{\mathrm{post}}$.

Two details the notation would otherwise hide. $P_k$ right-shifts by $k-1$ with a
zero pad \emph{before} averaging, so coarse token $j$ is the mean over original
indices $jk-1,\dots,jk+k-2$ --- not the plain group mean, and the first coarse token
is dragged toward zero by the pad. And $U_k$ is \texttt{repeat\_interleave}, which
is \emph{not} the adjoint of $P_k$ (that would carry a $1/k$ and the matching
shift); it is written $U_k$ rather than $P_k^{\dagger}$ for that reason.

Setting all three blocks to the \emph{same} shared module gives the weight-shared
hourglass: one block's parameters (204,373, identical to the flat 1-layer model,
against 600,917 unshared) with the coarse block's token reduction retained. Note
that every hourglass class ignores \texttt{n\_layers} and is always this 3-block
scaffold, so a flat control must be built at 3 blocks, not 1.

\section{Empirical anchors}

\begin{center}\small
\begin{tabular}{p{2.6cm}p{2.1cm}p{3.1cm}p{5.6cm}}
\toprule
mechanism & sets & parameter cost & measured effect \\
\midrule
Path integration & phase $\theta$ & $+448$ on $199{,}042$ ($+0.23\%$) & $+0.078$ parity ($16/16$); $+0.461$ torus \\[2pt]
PoPE             & magnitude $\mu$ & $\hd/2 \to \hd$ widening, $+320$ & $+0.027$ \emph{with} path int. Without it, both arms are on the floor --- see note \\[2pt]
Recursion        & block reuse & $0$ & $+0.056$ parity ($16/16$); matches or beats $3\times$ params \\[2pt]
Hierarchy        & token count & $0$ if weight-shared & length-dependent: $\approx 0$ at $L=16$, $+0.060$ at $L=512$ ($23/24$ seeds); $-22.8\%$ time \emph{at $L=2048$}, but see note (iv) \\[2pt]
InEKF            & $\theta \to \hat\theta$ & $+49{,}600$ on $204{,}373$ ($+24\%$) & no individual \emph{component} is load-bearing; the whole filter is $+0.06$ to $+0.13$, at OOD length only \\
\bottomrule
\end{tabular}
\end{center}

\noindent \textbf{Reading the table.} Parity figures are at $L=128$ trained at
$L=16$, from one batch of $8$ arms $\times$ $16$ seeds. Torus figures are the
position-axis factor average over both encodings ($+0.4379$ RoPE, $+0.4851$ PoPE).
\textbf{Three qualifications the columns cannot carry.}
(i) The hierarchy row is \emph{not} a parity-at-$L=128$ figure: its accuracy comes
from runs trained at $L=16$ and $L=512$, and its time saving is measured at
$L=2048$ --- at $L=128$ the same benchmark shows $+0.4\%$, i.e.\ no saving at all.
(ii) The InEKF's parameter cost is measured on a $204{,}373$-parameter torus model
and the path-integration cost on a $199{,}042$-parameter parity model; the two
percentages are not commensurable, and the InEKF has never been run on parity.
(iii) PoPE without path integration is a difference between two models sitting on
the measured floor of $0.506$ ($0.509$ vs $0.530$), with a seed spread wider than
the gap --- it is a floor artifact, not an effect, and is omitted from the column
for that reason.
(iv) The $-22.8\%$ is measured on \textsc{LoopedHourglass} against the unshared
flat scaffold, so it \emph{bundles} the shared block's parameter saving with the
pooling's compute saving. Isolated on text, hierarchy alone is $1.23\times$
throughput ($\approx -18.7\%$) at $-14.1\%$ memory. Attributing $-22.8\%$ to the
hierarchy axis alone overstates it.

\end{document}
